\chapter{Programy w C na Twoim procesorze}
\label{ch:c}

\metryczka{uruchomić skompilowany kod C na procesorze, który sam zbudowałeś, bez
systemu operacyjnego i bez libc. To terytorium, które znasz z pracy, tylko
tym razem pod spodem jest Twój własny sprzęt.}{1--2 wieczory}{07-c-bare-metal/}

\section{Toolchain}

Makefile wykrywa toolchain sam. Wystarczy jedno z dwóch:
\begin{sh}
sudo apt install lld                      # clang już jest, brakuje tylko linkera
sudo apt install gcc-riscv64-unknown-elf  # albo pełny toolchain GNU (umie też rv32)
\end{sh}
\code{make check-tools} pokazuje, czego brakuje. Moduł został sprawdzony z
clang~18 i \code{ld.lld}~18: wszystkie programy przechodzą na emulatorze, rdzeniu
jednocyklowym i potoku. Ścieżka GCC jest w Makefile, ale nie była testowana.

\section{Od C do pliku \texttt{.hex}}

\begin{figure}[h]
\centering
\begin{tikzpicture}[node distance=17mm, every node/.style={font=\sffamily\small}]
  \node[blok, fill=white] (c) {\code{hello.c}\\\code{rt.c}, \code{muldiv.c}};
  \node[blok, fill=white, below=4mm of c] (s) {\code{crt0.S}};
  \node[blok, right=20mm of $(c)!0.5!(s)$] (o) {\code{*.o}};
  \node[blok, right=of o] (elf) {\code{.elf}};
  \node[blok, right=of elf] (bin) {\code{.bin}};
  \node[blok, right=of bin, fill=zadaniejasny] (hex) {\code{.hex}};
  \node[blok, fill=white, above=6mm of elf] (ld) {\code{link.ld}};
  \draw[->, thick] (c) -- (o);
  \node[sygnal, above=1mm of o] {clang -c};
  \draw[->, thick] (s) -- (o);
  \draw[->, thick] (o) -- node[above, sygnal]{ld.lld} (elf);
  \draw[->, thick] (ld) -- (elf);
  \draw[->, thick] (elf) -- node[above, sygnal]{objcopy} (bin);
  \draw[->, thick] (bin) -- node[above, sygnal]{bin2hex} (hex);
\end{tikzpicture}
\caption{Kompilacja programu na nasz procesor.}
\end{figure}

Najważniejsze flagi:
\begin{itemize}
  \item \code{-march=rv32i -mabi=ilp32}: tylko instrukcje bazowe;
    \code{int}, \code{long} i wskaźnik mają po 32 bity,
  \item \code{-ffreestanding -nostdlib}: nie ma libc ani standardowego kodu startowego,
  \item \code{-mno-relax}: linker nie zamienia adresowania na względne
    wobec \code{gp} (nasz crt0 nie ustawia \code{gp}),
  \item \code{-O2}: warto zobaczyć, jak dobry kod generuje kompilator
    (\code{make dis P=bench}).
\end{itemize}

\section{Co się dzieje przed \texttt{main()}}

Procesor po resecie zaczyna od \code{pc = 0} i nic nie wie o C. Kod startowy
\plik{crt0.S} musi:
\begin{asm}
_start:
    la   sp, _stack_top        # 1. stos: C go potrzebuje do zmiennych lokalnych
    la   t0, __bss_start       # 2. zerowanie .bss: "static int x;" ma być 0
    la   t1, __bss_end
1:  bgeu t0, t1, 2f
    sw   zero, 0(t0)
    addi t0, t0, 4
    j    1b
2:  call main                  # 3. program
    ...                        # 4. a0 == 0 ? TOHOST = 1 : TOHOST = (a0 << 1) | 1
\end{asm}

Sekcji \code{.data} nie trzeba kopiować, bo cały obraz programu ląduje w RAM przez
\code{\$readmemh}. Na prawdziwym mikrokontrolerze \code{.data} leży we flashu i
crt0 kopiuje ją do RAM.

\section{Skrypt linkera}

\begin{figure}[h]
\centering
\begin{tikzpicture}[x=1cm, y=1cm]
  \foreach \y/\h/\n/\c in {0/0.9/{.text (crt0 najpierw!)}/akcentjasny, 0.9/0.6/.rodata/akcentjasny,
                         1.5/0.5/.data/zadaniejasny, 2.0/0.5/.bss/zadaniejasny,
                         2.5/1.6/{wolne}/white, 4.1/0.9/{stos $\downarrow$}/pulapkajasny} {
    \draw[thick, fill=\c] (0,-\y) rectangle (5,-\y-\h);
    \node[font=\sffamily\small] at (2.5,-\y-\h/2) {\n};
  }
  \node[sygnal, anchor=east] at (-0.1,0) {0x0000};
  \node[sygnal, anchor=east] at (-0.1,-2.0) {\_\_bss\_start};
  \node[sygnal, anchor=east] at (-0.1,-2.5) {\_\_bss\_end};
  \node[sygnal, anchor=east] at (-0.1,-5.0) {0x10000 = \_stack\_top};
\end{tikzpicture}
\caption{Układ pamięci programu (64 KiB RAM).}
\end{figure}

Symbole \code{\_\_bss\_start}, \code{\_\_bss\_end} i \code{\_stack\_top} definiuje
\plik{link.ld}, a używa ich \plik{crt0.S}. W ten sposób linker przekazuje adresy do
kodu.

\section{Brakujące instrukcje: \texttt{\_\_mulsi3} i spółka}

RV32I nie ma mnożenia ani dzielenia. Gdy piszesz \code{a * b}, kompilator wstawia
\code{call \_\_mulsi3}. Zwykle dostarcza ją \code{libgcc} albo
\code{compiler-rt}, a u nas piszesz ją sam (zadanie 7.2). Zasady:

\begin{itemize}
  \item w \plik{muldiv.c} \textbf{nie wolno} używać \code{*}, \code{/} ani
    \code{\%}, bo kompilator wywołałby te same funkcje i powstałaby nieskończona
    rekurencja,
  \item dzielenie przez zero ma dawać to samo co instrukcje RISC-V M:
    $x / 0 = $ \code{0xFFFFFFFF}, $x \bmod 0 = x$,
  \item dzielenie ze znakiem zaokrągla do zera: $-7/2 = -3$, $-7 \bmod 2 = -1$.
\end{itemize}

Mnożenie podaję jako gotowy przykład, a dzielenie piszesz sam. Mnożenie „pisemne” w systemie dwójkowym:
\begin{cc}
uint32_t __mulsi3(uint32_t a, uint32_t b) {
    uint32_t r = 0;
    while (b) {
        if (b & 1) r += a;     // bit mnożnika ustawiony: dodaj przesuniętą mnożną
        a <<= 1;
        b >>= 1;
    }
    return r;
}
\end{cc}
Dzielenie robi się analogicznie (\emph{restoring division}): przesuwaj resztę w
lewo, dosuwaj kolejny bit dzielnej, a gdy reszta $\geq$ dzielnik, odejmij i ustaw
bit ilorazu.

Kompilator potrafi też sam wstawić wywołania \code{memset} i \code{memcpy} (np.
przy zerowaniu tablicy na stosie), nawet jeśli ich nie używasz. Dlatego są w
\plik{lib/rt.c}.

\section{MMIO i \texttt{volatile}}

\begin{cc}
#define MMIO_PUTCHAR ((volatile uint32_t *)0x10000004)
*MMIO_PUTCHAR = 'A';
\end{cc}
Bez \code{volatile} kompilator mógłby uznać dwa zapisy pod ten sam adres za zbędne
i usunąć jeden. Na tym poziomie nie ma różnicy między zmienną a rejestrem
urządzenia: to ten sam \code{sw}, tylko pod inny adres, który testbench (albo
prawdziwy dekoder adresów) kieruje do urządzenia.

\section{Konwencja wywołań (ilp32)}

Argumenty trafiają do \code{a0}--\code{a7}, wynik do \code{a0}, adres powrotu do
\code{ra}. Rejestry \code{s0}--\code{s11} funkcja musi zachować, więc zapisuje je
na stosie w prologu. \code{t*} i \code{a*} może niszczyć. Stos jest wyrównany do
16 bajtów. Typowy prolog i epilog:
\begin{asm}
    addi sp, sp, -16
    sw   ra, 12(sp)
    sw   s0, 8(sp)
    ...
    lw   s0, 8(sp)
    lw   ra, 12(sp)
    addi sp, sp, 16
    ret
\end{asm}

\section{Zadania}

\begin{zadanie}{7.1 --- Hello, world}
\code{make} (kompilacja), \code{make emu P=hello} (emulator), \code{make run P=hello}
(Twój rdzeń; \code{CORE=08} uruchamia potok). Pod jakim adresem leży
\code{counter} z \plik{hello.c} i w której jest sekcji? Co by się stało bez
zerowania \code{.bss}? Ile instrukcji wykonuje \code{hello}, a ile z nich to crt0?
\end{zadanie}

\begin{zadanie}{7.2 --- \code{lib/muldiv.c}}
\code{make emu P=muldiv\_test}, potem \code{make run P=muldiv\_test}. Wektory
testowe obejmują dzielenie przez zero, \code{INT\_MIN / -1} i liczby ujemne.
\end{zadanie}

\begin{zadanie}{7.3 --- Benchmark}
\code{make run P=bench} wypisuje wyniki i liczbę cykli dla CRC32, sita
Eratostenesa, sortowania i mnożenia macierzy. Która część jest najdroższa?
Sprofiluj ją na poziomie sprzętu:
\begin{sh}
cut -d' ' -f1 build/bench.trace | sort | uniq -c | sort -rn | head
\end{sh}
i znajdź najczęstsze adresy w \plik{build/bench.dis}.
\end{zadanie}

\begin{zadanie}{7.4 --- Własny program}
Napisz \plik{prog/coś.c}: liczby pierwsze do 10\,000, drzewo BST w statycznej
tablicy albo „grę w życie” wypisywaną znakami. \code{main} zwraca 0 jako PASS.
\end{zadanie}

\begin{gotowe}
\code{make test} → wszystkie PASS na emulatorze i na rdzeniu.
\end{gotowe}

\begin{kontrolne}
  \item Co musi zrobić crt0, zanim wywoła \code{main}? Czego nie musi robić u nas, a musi na mikrokontrolerze?
  \item Skąd crt0 zna adresy początku i końca \code{.bss}?
  \item Dlaczego w \plik{muldiv.c} nie można użyć operatora \code{*}?
  \item Co się stanie, gdy usuniesz \code{volatile} z definicji \code{MMIO\_PUTCHAR} i skompilujesz z \code{-O2}?
\end{kontrolne}
